\section{Background}
\label{sec:rs-background}

The Recursive Spawning Protocol (RSP) operates at the intersection of several research traditions: agent lifecycle management in multi-agent AI systems, accountability and provenance in distributed computing, fixed versus mutable delegation structures, hierarchical task decomposition, and immutable parent-chain architectures. This section situates RSP within each tradition, traces the conceptual gaps that motivate it, and establishes the BX3 Framework as the minimal sufficient substrate for RSP's accountability guarantees.

% -------------------------------------------------------
\subsection{Agent Lifecycle Management in Multi-Agent Systems}
% -------------------------------------------------------

Multi-agent systems (MAS) research has long recognized that the lifecycle of an agent --- its creation, activation, operation, suspension, and termination --- is a first-class design concern, not merely an implementation detail. The foundational taxonomy introduced by Stone and Veloso \cite{stone2000} classifies multi-agent systems along axes of agent heterogeneity, agent knowledge quality, environment accessibility, and the nature of agent interactions. Within this taxonomy, recursive spawning is most relevant to cooperative, homogeneous systems operating in partially accessible environments --- precisely the regime in which modern agentic runtimes such as AutoGPT \cite{autoGPT2023}, LangChain \cite{langchain2023}, and AutoGen \cite{autogen2023} operate.

Classical agent lifecycle management was primarily an engineering concern: spawning, scheduling, and terminating agents to maximize throughput or minimize latency, with accountability residing with the system operator who configured the agent population. As AI agents move into healthcare, infrastructure control, and financial decision-making, this assumption becomes untenable. When a spawned agent's action produces real-world consequences, the question of \emph{who authorized the spawn}, through what chain of delegation, and with what scope constraints is no longer academic --- it is a legal, regulatory, and safety question with no clean answer in conventional architectures.

Ferber and M\"uller \cite{ferber1996} formalized agent lifecycle interactions as interactions between an agent's lifecycle state machine and the organizational context in which it operates, demonstrating that organizational roles --- not just individual agent states --- must be modeled to ensure system-wide coherence. Padgham and Thiel \cite{padgham2001} further established that incomplete lifecycle management, particularly the failure to distinguish between suspension and termination states, introduces subtle but systemic failures: orphaned or zombie agents that continue to consume resources and emit effects after their parent has exited.

More recent work has surfaced the specific problem of \emph{unbounded spawning}. Zhuge et al. \cite{zhuge2024} document the agent chain problem: LLM-based agents that recursively spawn children to handle sub-tasks, without a defined termination condition or depth ceiling, can produce chains of dozens of intermediate agents whose aggregate authority differs substantially from what any single agent was authorized to exercise. Without a structural enforcement mechanism, the delegation chain grows organically and the system's effective authority expands beyond any human author's intent.

Recent protocols have begun addressing this gap directly. The Human Delegation Provenance (HDP) Protocol \cite{hdp-draft,hdp-draft-2} identifies the core vulnerability in multi-hop delegations: in human $\rightarrow$ orchestrator $\rightarrow$ sub-agent $\rightarrow$ tool-agent chains, verifying that terminal actions were genuinely authorized by a human requires reconstructing the full delegation chain, a task that OAuth 2.0, JWTs, and UCAN do not adequately address for append-only, human-provenance requirements in agentic systems. HDP proposes cryptographic tokens capturing each delegation hop as a signed entry in an append-only chain, enabling offline verification using only the issuer's public key and session identifier.

The Agent Lifecycle Protocol (ALP) \cite{alp2025} extends this by defining canonical lifecycle events --- Genesis, Fork, Migration, Retirement, Succession, and Decommission --- with explicit state-machine semantics, transition rules, and boundary hooks, alongside a genealogical registry capturing both genetic (model, architecture) and epigenetic (configuration, memory, reputation history) lineage. PROV-AGENT \cite{prov-agent2025} further demonstrates this convergence, introducing a provenance model extending W3C PROV \cite{prov-standard} with the Model Context Protocol to link agent decisions and prompts with end-to-end workflow data. Unified Agent Lifecycle Management (UALM) \cite{ualm2025} provides a complementary five-layer blueprint for end-to-end governance of agent fleets: identity and persona registry, orchestration and cross-domain mediation, context-bounded memory, runtime policy enforcement with kill-switch triggers, and lifecycle management linked to credential revocation and audit logging.

Together, these protocols demonstrate that the field has identified the accountability gap in agentic lifecycles and is converging on cryptographic, append-only delegation records as the solution. RSP's contribution is providing the structural enforcement layer --- through the immutable parent pointer --- that makes these delegation records immutable by construction rather than conventionally immutable by policy. The core problem RSP addresses is \emph{spawn chain opacity}: the inability to determine, at any point during or after execution, whether a given agent's authority traces to a human principal or to an autonomous predecessor.

% -------------------------------------------------------
\subsection{Accountability and Provenance in Distributed Systems}
% -------------------------------------------------------

Accountability in distributed systems has historically been analyzed through the lens of audit logs, non-repudiation, and cryptographic provenance. The W3C PROV data model \cite{prov-standard} provides the foundational vocabulary: entities (the things being accounted for), activities (the actions performed), and agents (the principals responsible). Provenance records capture the chain of custody for data and decisions, enabling post-hoc verification that a given output was generated by a specific process under specific conditions. In distributed ledgers such as Bitcoin \cite{nakamoto2008bitcoin} and Hyperledger Fabric \cite{hyperledger2015}, this principle is operationalized as an append-only, hash-chained transaction log: the history cannot be rewritten without breaking the hash chain, and any breakage is immediately detectable.

The extension of this principle to multi-agent AI systems raises challenges absent from classical distributed systems. Agents are not static processes --- they are autonomous entities that can modify their behavior mid-execution, inherit and propagate authority across spawn chains, and take actions whose consequences may not be observable until long after the chain has grown. The problem of accountability in these systems is therefore not merely the problem of recording what happened, but of recording it in a way that is \emph{tamper-evident with respect to the delegation chain specifically}: the record must make it possible to determine, for any agent $a$, which human principal authorized the chain of decisions that led to $a$'s current state.

This requirement motivates what we term the \emph{immutability of authorization} requirement. The provenance record must capture not only data transformations but the \emph{delegation chain} itself as a first-class artifact. If the delegation chain is mutable --- if parent pointers can be updated retroactively --- then any accountability guarantee derived from the provenance record is conditional on the absence of a retroactive edit. Certificate Transparency \cite{laurie2013certificate} illustrates this principle in the domain of TLS certificates: append-only logging combined with a verification mechanism (the Merkle audit path) makes certificate mis-issuance detectable after the fact. The same architectural pattern applies to agent delegation chains: an immutable, append-only log of spawn events makes unauthorized delegation detectable.

Systems such as CollectiveOS \cite{collectiveos} and the Proof-of-Behavior Protocol \cite{pob-draft} demonstrate the emerging recognition that accountability guarantees in autonomous AI systems are only as strong as the immutability guarantees of the provenance record itself. Recent work on immutable audit trails for Non-Human Identities (NHIs) \cite{nhi-lineage} further establishes that the NHI problem --- the challenge of establishing persistent, verifiable identity for autonomous agents --- is coextensive with the delegation chain immutability problem: an agent's identity is, in the relevant sense, its chain of authorizations.

% -------------------------------------------------------
\subsection{Fixed Versus Mutable Delegation Chains}
% -------------------------------------------------------

Delegation chains in multi-agent systems take two structurally different forms with profoundly different accountability properties. In a \emph{fixed delegation chain}, the parent-child relationship is established at spawn time and is immutable thereafter: the child's parent pointer cannot be changed, the child's authority cannot be re-delegated except through the defined spawn protocol, and the chain terminates at a predetermined human anchor. In a \emph{mutable delegation chain}, the parent-child relationship is treated as runtime state that can be updated, redirected, or reassigned as conditions change.

Mutable delegation chains offer flexibility: a child agent can be reparented to a different parent if the organizational structure changes, if a parent fails, or if load balancing requires reassignment. This flexibility comes at a significant accountability cost. If the parent-child relationship is mutable at runtime, any single action by a child agent may have been authorized by one parent at one moment and a different parent at another --- and the effective authorization for any given action may not be the authorization that was recorded at spawn time. The chain of provenance becomes conditional on the state of a mutable mapping rather than fixed by an immutable record. This introduces the \emph{delegatee ambiguity problem}: after an event, it may be impossible to determine which principal effectively authorized the action, because the delegation chain has been rewritten since the action occurred.

A fixed delegation chain sacrifices this flexibility for a stronger accountability guarantee. Each child's authority is traceable to a single, immutable parent pointer that chains unbroken back to a human anchor. The Fact Layer enforces immutability of the parent pointer through its write protocol, rejecting any request to overwrite $\alpha(n)$ after provisioning regardless of who issued the request. The immutability is structural, not contractual: any attempt to overwrite $\alpha(n)$ is rejected at the protocol level, before execution. This contrasts with ACL-based immutability in conventional systems, where administrative compromise leads directly to immutability violation; BX3's protocol-level enforcement has no administrative override path.

The tradeoff corresponds to what the distributed systems literature terms \emph{strong consistency} versus \emph{eventual consistency} \cite{vogels2009eventual}: fixed chains correspond to strong consistency (every read of the delegation chain returns the same value, established at spawn time), while mutable chains correspond to eventual consistency (the current parent pointer is eventually correct, but intermediate states may be ambiguous). For accountability purposes, strong consistency is the required property: the authorization for any agent action must be deterministically recoverable at any time, not just eventually.

SentinelAgent \cite{sentinelagent} formalizes the secure delegation chain requirement specifically for federal multi-agent AI systems, demonstrating that the combination of cryptographic chain-of-custody, formal verification of delegation constraints, and an immutable parent pointer produces the strongest available guarantee against unauthorized delegation propagation. RSP adopts this principle and operationalizes it within the BX3 Framework's three-layer architecture.

% -------------------------------------------------------
\subsection{Hierarchical Task Decomposition in AI}
% -------------------------------------------------------

Hierarchical task decomposition (HTD) --- the recursive decomposition of a complex task into simpler subtasks --- is a foundational technique in classical AI planning. Classical HTN planning \cite{ghallab2004htn,erpl} decomposes a high-level task into a hierarchy of sub-tasks according to predefined method schemas, where each decomposition is a planning decision made \emph{before execution}. The decomposition tree is a planning artifact: it represents the planner's reasoning about how to accomplish the goal, not a runtime hierarchy of autonomous entities. Subtasks in classical HTN are subroutine calls; they do not have independent lifecycle states, independent authority chains, or the capacity to spawn further sub-agents.

Modern LLM-based agents have adopted a fundamentally different form of HTD. Rather than decomposing a task into sub-tasks according to a predefined schema, an LLM agent can observe that a current action has failed, infer that the failure stems from task complexity beyond its provisioned scope, and autonomously spawn a child agent to handle a specific sub-task \cite{press2022,yao2022,thread2025}. The decomposition is reactive and context-sensitive: the spawn tree grows dynamically based on the agent's real-time observations. THREAD \cite{thread2025} provides perhaps the closest conceptual relative to RSP, treating LLM generation as a dynamic multi-threaded process where a parent thread can spawn child threads to handle sub-tasks and pause until those children return results. This dynamic recursive spawning enables hierarchical task decomposition while managing context length, achieving state-of-the-art performance on agent-task and data-grounded QA benchmarks with smaller models.

What connects classical HTN and modern LLM-based decomposition is the spawn tree: a hierarchical structure in which each node represents a subtask and edges represent the decomposition relationship. In both traditions, the tree is typically mutable and implicit --- decomposition decisions are not recorded as first-class authorization artifacts, and child nodes are not distinguished from parent nodes in terms of their accountability status. A classical HTN planner's subtasks are subroutine calls without independent accountability chains; an LLM agent's spawned sub-agents are independent processes, but the chain of authorization is implicit in the agent's context window rather than explicit in an immutable warrant structure.

Dynamic HTD in open-loop agentic systems has been identified as a vector for \emph{autonomous drift}: the progressive expansion of agent behavior beyond authorized scope, precisely because decomposition is unconstrained \cite{autoGPT2023}. The spawn necessity threshold --- the Purpose Layer's requirement that each child justify why it cannot be replaced by the parent --- addresses a specific HTD failure mode: \emph{artificial fragmentation}, where a parent spawns children not to exploit genuine parallelism or specialized capability but to circumvent context-window limits or create the appearance of multi-agent orchestration. This mirrors \emph{serialization of parallelizable tasks} in classical planning --- the decomposition of a task that should execute concurrently into a sequential chain consuming more depth without delivering more capability. The Purpose Layer's spawn necessity check forces the parent to articulate why spawning is necessary, making artificial fragmentation auditable and rejectable.

Recursive Spawning adopts the hierarchical decomposition paradigm from both traditions but adds the critical accountability layer absent from both: each node in the decomposition tree is a first-class accountability entity with an immutable parent pointer, a defined termination condition, and a cryptographically verifiable chain back to a human anchor. Decomposition is not merely a planning technique; it is an authorization event that creates a new accountable entity. The decomposition tree is therefore not just a planning artifact but an organizational chart of delegated authority.

% -------------------------------------------------------
\subsection{The BX3 Framework as Substrate for Immutable Parent Chains}
% -------------------------------------------------------

The BX3 Framework's three-layer model --- Purpose Layer, Bounds Engine, and Fact Layer --- provides the minimal sufficient architectural substrate for RSP's accountability guarantees \cite{bx3framework2026}. The mapping between RSP components and BX3 layers is not a design choice among alternatives; it is a structural necessity. Each RSP guarantee requires exactly the capabilities that one BX3 layer provides and no other layer can provide.

The \emph{immutable parent pointer} is a Fact Layer write protocol constraint. The Fact Layer is the only layer in the BX3 architecture capable of producing non-reversible physical effects. The parent pointer $\alpha(n)$ is established once, at node provisioning, through a Fact Layer write operation that is atomic and immediately hash-chained to the preceding ledger entry. After this write, no actor --- including the Purpose Layer that originated it --- can modify $\alpha(n)$. This immutability is not a policy enforced by convention; it is a structural constraint enforced by the Fact Layer's write validation. Any attempt to overwrite $\alpha(n)$ is rejected at the protocol level, before execution. This contrasts with ACL-based immutability in conventional systems: administrative compromise in ACL-based systems leads directly to immutability violation, whereas BX3's protocol-level enforcement has no administrative override path.

The \emph{depth bound} $D_{\max}$ is a Purpose-Layer-defined, Bounds-Engine-enforced, Fact-Layer-hardened constraint. The Purpose Layer defines the maximum permissible chain depth as part of the spawn authorization policy. The Bounds Engine evaluates each proposed spawn against this limit, returning \texttt{SPAWN\_REFUSED} if the spawn would exceed it. The Fact Layer's pre-commit hook enforces this limit structurally: any spawn operation that would produce $|C| \geq D_{\max}$ is rejected before execution, regardless of the Bounds Engine's return. The constraint is thus enforced three times --- by the layer that defines it, the layer that evaluates against it, and the layer that structurally blocks violations.

The \emph{forensic ledger} is a Fact Layer artifact. Every spawn event, state transition, Purpose Layer directive, and exception condition is recorded in the append-only ledger maintained by the Fact Layer. Each entry is hash-chained to its predecessor, making the ledger tamper-evident: any insertion, deletion, or modification of a ledger entry breaks the hash chain and is detectable by any observer with read access. The ledger provides the evidentiary foundation for RSP's correctness properties --- drift prevention, chain integrity, and termination --- because it contains a complete, immutable record of every decision point in the chain's lifecycle.

The \emph{human anchor} $h(n)$ is a Purpose Layer assignment. At provisioning, the Purpose Layer records the identity of the human principal who authorized the root spawn event. This assignment is immutable for the lifetime of the chain: for all nodes $n$ in chain $C$, $h(n) = h(n_0)$ regardless of chain depth. This is the architectural expression of the upstream BX3 accountability guarantee: systems built on the BX3 Framework fail upward into human consciousness, never downward into autonomous chaos.

The structural necessity of the three-layer mapping is established formally in Section~\ref{sec:rs-integration}. Removing any layer causes the corresponding guarantee to fail. Without the Purpose Layer, there is no authoritative human anchor and no source of spawn authorization policy. Without the Bounds Engine, depth limits become unenforceable conventions subject to override by operational urgency. Without the Fact Layer, immutability of the parent pointer and tamper-evidence of the ledger are promises rather than constraints. The three-layer architecture is not a design pattern --- it is the minimal structure that makes RSP's correctness properties unconditionally guaranteed rather than conditionally expected policies.
